\documentclass[11pt]{article}
\usepackage[T1]{fontenc}
\usepackage{amsmath,amssymb,amsthm,mathtools}
\usepackage{graphicx}
\usepackage[margin=0.88in]{geometry}
\usepackage{microtype}
\usepackage[hidelinks]{hyperref}

\newtheorem{theorem}{Theorem}
\newtheorem{lemma}[theorem]{Lemma}
\newtheorem{remark}[theorem]{Remark}
\newcommand{\cL}{\mathcal L}

\title{Positive linearization of finite-dimensional isotone contractions\\
\large The finite-dimensional case of Remark 4.9(c) in arXiv:1810.09684}
\author{Autonomous proof packet --- agent\_lane\_12}
\date{11 August 2026}

\begin{document}
\maketitle

\begin{abstract}
Chill, Fiorenza, and Kr\'ol ask whether exact interpolation for positive
linear operators always implies exact interpolation for order-preserving
Lipschitz operators with solid lattice domain.  We prove the implication for
every finite-dimensional compatible Banach-lattice couple.  The key is a
pairwise linearization theorem: for any two points $f,h$ in the domain of an
order-preserving endpoint contraction $S$, there is one positive admissible
linear endpoint contraction $T$ such that
$T(f-h)=S(f)-S(h)$.  We smooth $S$ inside the order box generated by $f,h$;
its derivatives are positive endpoint contractions, and their average along
the segment gives $T$ in the limit.  The argument includes proper solid
domains, the feature not covered by the source's observation about
everywhere-defined Lipschitz maps.  The general infinite-dimensional
question remains open.
\end{abstract}

\section{Source problem and result}

The source is R.~Chill, A.~Fiorenza, and S.~Kr\'ol,
\emph{Interpolation of nonlinear positive or order preserving operators on
Banach lattices}, arXiv:1810.09684~\cite{source}.  Theorem 4.8 proves
\begin{align*}
 \text{exact partially $K$-monotone}
 &\Longleftrightarrow \text{exact positive Gagliardo--Peetre}\\
 &\Longrightarrow \text{exact order-preserving Lipschitz}
 \Longrightarrow \text{exact positive linear}.
\end{align*}
Remark 4.9(c), on source PDF page 18, asks whether the last implication can
always be reversed.

\begin{center}
\includegraphics[width=0.97\textwidth]{figures/open_problem_crop.png}
\end{center}

An interpolation couple $\vec E=(E_0,E_1)$ is understood in the source's
sense: $E_0,E_1$, and $E_0+E_1$ are Banach lattices, and the endpoint spaces
are order ideals in the sum.  For an admissible linear operator set
\[
 \|T\|_{\cL(\vec E)}
 =\max_{i=0,1}\|T:E_i\to E_i\|.
\]

\begin{theorem}[Finite-dimensional affirmative answer]\label{thm:main}
Let $\vec E=(E_0,E_1)$ be a compatible Banach-lattice couple for which
$E:=E_0+E_1$ is finite-dimensional, and let $X$ be an intermediate space.
If $X$ is an exact interpolation space for positive linear operators on
$\vec E$, then it is an exact interpolation space for admissible
order-preserving Lipschitz operators whose domains are solid lattices.
Thus implication (iv)$\Rightarrow$(iii) of source Theorem 4.8 holds for every
finite-dimensional compatible Banach-lattice couple.
\end{theorem}

The theorem follows immediately from the stronger pairwise statement below.

\section{Pairwise positive linearization}

\begin{lemma}\label{lem:linearize}
Let $E=E_0+E_1$ be finite-dimensional.  Suppose that
$S:D\to E$ is order preserving, $D\subset E$ is a solid vector lattice, and
for some $L\ge0$,
\begin{equation}\label{eq:lip}
 \|Sx-Sy\|_{E_i}\le L\|x-y\|_{E_i}
 \quad\text{whenever }x,y\in D\text{ and }x-y\in E_i,
 \qquad i=0,1.
\end{equation}
Then for every $f,h\in D$ there exists a positive admissible linear operator
$T:E\to E$ such that
\[
 \|T\|_{\cL(\vec E)}\le L,
 \qquad T(f-h)=S(f)-S(h).
\]
\end{lemma}

\begin{proof}
The case $f=h$ is immediate.  A finite-dimensional vector lattice is
lattice-isomorphic to coordinate-ordered $\mathbb R^m$, so we use that model.
Put
\[
 a=f\wedge h,\qquad b=f\vee h.
\]
Since $D$ is a lattice, $a,b\in D$; since it is solid, the complete order box
$[a,b]$ lies in $D$.  Let $B$ be the coordinate band spanned by the active
coordinates $j$ for which $a_j<b_j$, and let $P$ be its band projection.
Every band projection is positive and contractive for every lattice norm.
Moreover $P(f-h)=f-h$.

We first note that the restriction of $S$ to the relative interior
\[
 U=\{x\in a+B: a_j<x_j<b_j\text{ on every active coordinate}\}
\]
is locally Lipschitz in an ordinary finite-dimensional norm.  Indeed, the
order ideals $E_0,E_1$ are coordinate subspaces and their sum is $E$.
Choose disjoint coordinate band projections $R_0,R_1$ with
$R_0+R_1=I$ and $R_iE\subset E_i$.  If $x,y$ lie in an inner order box, then
the coordinatewise mixed point $z=x+R_0(y-x)$ lies in the same box, and
\begin{align*}
 S(y)-S(x)&=[S(y)-S(z)]+[S(z)-S(x)].
\end{align*}
The two differences are controlled by \eqref{eq:lip} in $E_1$ and $E_0$,
respectively.  This proves the asserted local continuity and permits
ordinary convolution inside $U$.

Let $\rho_\delta$ be a nonnegative smooth mollifier on $B$, supported in a
ball of radius $\delta$ and of integral one.  On an inner box define
\[
 S_\delta(x)=\int_B S(x-z)\rho_\delta(z)\,dz.
\]
This map is smooth.  It remains order preserving: whenever $x\le y$ and all
translates stay in $U$, each integrand in
$S(y-z)-S(x-z)$ is positive.  Consequently
$DS_\delta(x):B\to E$ is a positive linear map.

Convolution also preserves both endpoint Lipschitz estimates.  If
$v\in B\cap E_i$, then
\begin{align*}
 \|S_\delta(x+tv)-S_\delta(x)\|_{E_i}
 &\le \int_B
 \|S(x+tv-z)-S(x-z)\|_{E_i}\rho_\delta(z)\,dz\\
 &\le L|t|\,\|v\|_{E_i}.
\end{align*}
After division by $|t|$ and passage to the limit, the operator
\[
 A_x:=DS_\delta(x)P:E\to E
\]
is positive, admissible, and satisfies
$\|A_x\|_{\cL(\vec E)}\le L$.

Let $c=(a+b)/2$ and, for $0<\varepsilon<1$, set
\[
 f_\varepsilon=(1-\varepsilon)f+\varepsilon c,
 \qquad
 h_\varepsilon=(1-\varepsilon)h+\varepsilon c.
\]
These points and the segment between them lie in an inner box.  Choose
$\delta$ small enough for the preceding construction and average the
derivatives along that segment:
\[
 T_{\varepsilon,\delta}
 =\int_0^1
 DS_\delta\bigl(h_\varepsilon+t(f_\varepsilon-h_\varepsilon)\bigr)P\,dt.
\]
It is again positive and
$\|T_{\varepsilon,\delta}\|_{\cL(\vec E)}\le L$.  The fundamental theorem
of calculus gives
\begin{equation}\label{eq:approx}
 T_{\varepsilon,\delta}(f-h)
 =\frac{S_\delta(f_\varepsilon)-S_\delta(h_\varepsilon)}{1-\varepsilon}.
\end{equation}

First let $\delta\downarrow0$ and then $\varepsilon\downarrow0$.  The local
continuity established above shows that the right side of
\eqref{eq:approx} converges to $S(f)-S(h)$.  The family of positive
admissible operators with endpoint norms at most $L$ is bounded in the
finite-dimensional space $\cL(E)$: the fixed decomposition
$E=R_0E+R_1E$ bounds its norm on all of $E$.  It is also closed.  Hence a
subsequence of $T_{\varepsilon,\delta}$ converges to a positive admissible
operator $T$ with $\|T\|_{\cL(\vec E)}\le L$, and \eqref{eq:approx} yields
$T(f-h)=S(f)-S(h)$.
\end{proof}

\begin{proof}[Proof of Theorem \ref{thm:main}]
Let $S$ be an order-preserving admissible Lipschitz operator with solid
lattice domain, let $L=\|S\|_{\mathrm{Lip}}$, and choose $f,h$ in its domain
with $f-h\in X$.  Lemma \ref{lem:linearize} supplies a positive admissible
linear $T$ satisfying $T(f-h)=S(f)-S(h)$ and
$\|T\|_{\cL(\vec E)}\le L$.  Exactness of $X$ for positive linear operators
therefore gives
\[
 \|S(f)-S(h)\|_X
 =\|T(f-h)\|_X
 \le L\|f-h\|_X,
\]
which is precisely source assertion (iii).
\end{proof}

\section{Scope and upgrade audit}

The proof covers arbitrary finite-dimensional compatible lattice norms and,
crucially, arbitrary solid lattice domains.  It is stronger than merely
extending $S$: the linear operator may depend on the selected pair, which is
all exact interpolation needs.

The finite-dimensional hypothesis enters twice.  Convolution occurs on a
Lebesgue-measured order box, and the bounded family of averaged derivatives
is compact in $\cL(E)$.  For a general Banach lattice, restricting to the
two-dimensional lattice spanned by $(f-h)^+$ and $(f-h)^-$ produces only a
positive operator on that sublattice.  Extending it to the whole couple while
preserving both endpoint constants is exactly the positive-extension/band
gluing property not available in the source's general setting.  Endpoint
Lipschitz assumptions also give no continuity in the $E_0+E_1$ topology for
differences lying in neither endpoint, so finite-band approximation does not
close the gap.  These are genuine obstructions to the present method, not
claims that the full implication is false.

Bounded searches on 11 August 2026 covered the run's cheap indexes and
arXiv/web queries by exact title, authors, Remark 4.9(c), implication
(iv)$\Rightarrow$(iii), and finite-dimensional order-preserving
linearization.  They found the source and related extension literature but
no later resolution or matching finite-dimensional theorem.  This is not a
priority claim.

\paragraph{Human-review recommendation.}
Accept as a substantial partial result.  Review chiefly the local
two-endpoint continuity argument, preservation of endpoint constants under
mollification, and the final compactness passage.  Do not record the general
infinite-dimensional problem as solved.

\begin{thebibliography}{9}
\bibitem{source}
R.~Chill, A.~Fiorenza, and S.~Kr\'ol,
\emph{Interpolation of nonlinear positive or order preserving operators on
Banach lattices}, arXiv:1810.09684 (2018); Positivity 24 (2020), 507--532,
Theorem 4.8 and Remark 4.9(c).
\end{thebibliography}

\end{document}
